%------------------------------------------------------------------------------%
\begin{question}
  Determine os autovalores, os autovetores e os autoespaços da matriz
  \[
    B =
    \begin{bmatrix}
      4 & 1 & 0 \\
      1 & 4 & 0 \\
      0 & 0 & 2
    \end{bmatrix}
  \]
\end{question}

%------------------------------------------------------------------------------%
\begin{resolution}{Autovalores}
  Polinômio característico
  \begin{align*}
    p(\lambda)
    & = \det(B-\lambda I) \\
    & = \det\left(
          \begin{bmatrix}
            2 & 1 & 0 \\
            0 & 2 & 0 \\
            0 & 0 & 5
          \end{bmatrix}
          -
          \begin{bmatrix}
            \lambda & 0       & 0       \\
            0       & \lambda & 0       \\
            0       & 0       & \lambda
          \end{bmatrix}
          \right)
    \;=\; \begin{vmatrix}
      2-\lambda & 1         & 0         \\
      0         & 2-\lambda & 0         \\
      0         & 0         & 5-\lambda
    \end{vmatrix}
    \\
    & = (2-\lambda)^2(5-\lambda)
  \end{align*}
  Logo, os autovalores são
  \[
    \lambda_{1} = 2
    \qquad
    \lambda_{2} = 5
  \]
\end{resolution}

%------------------------------------------------------------------------------%
\begin{resolution}{Autovetores de \(\lambda=2\)}
\begin{columns}
\column{0.5\textwidth}
  \begin{align*}
    (B-2I)\vec{v} & = \vec{0}
    \\[1mm]
    \begin{bmatrix}[c]
      2-2 & 1   & 0   \\
      0   & 2-2 & 0   \\
      0   & 0   & 5-2
    \end{bmatrix} & =
    \begin{bmatrix}
      0 \\
      0 \\
      0
    \end{bmatrix}
    \\[3mm]
    \begin{bmatrix}
      0 & 1 & 0 \\
      0 & 0 & 0 \\
      0 & 0 & 3
    \end{bmatrix}
    \begin{bmatrix}
      x \\
      y \\
      z
    \end{bmatrix}
    & =
    \begin{bmatrix}
      0 \\
      0 \\
      0
    \end{bmatrix}
   \end{align*}
\column{0.4\textwidth}
  Solução
  \begin{align*}
    x & = t \in \R \\
    y & = 0 \\
    z & = 0
  \end{align*}
  Um autovetor
  \[
    \vec{v}_1
    = \vetortri{1}{0}{0}
  \]
\end{columns}
\end{resolution}

%------------------------------------------------------------------------------%
\begin{resolution}{Autoespaço de \(\lambda=2\)}
  \[
    E(2)
    = \left\{
      \begin{bmatrix}
        t \\
        0 \\
        0
      \end{bmatrix}
      \st t \in \R
    \right\}
    = \ger\left\{
      \begin{bmatrix}
        1 \\
        0 \\
        0
      \end{bmatrix}
    \right\}
  \]
\end{resolution}

%------------------------------------------------------------------------------%
\begin{resolution}{Autovetores de \(\lambda=5\)}
\begin{columns}
\column{0.5\textwidth}
  \begin{align*}
    (B-5I)\vec{v} & = \vec{0}
    \\[1mm]
    \begin{bmatrix}[c]
      2-5 & 1   & 0   \\
      0   & 2-5 & 0   \\
      0   & 0   & 5-5
    \end{bmatrix} & =
    \begin{bmatrix}
      0 \\
      0 \\
      0
    \end{bmatrix}
    \\[3mm]
    \begin{bmatrix}
      -3 & 1  & 0 \\
      0  & -3 & 0 \\
      0  & 0  & 0
    \end{bmatrix}
    \begin{bmatrix}
      x \\
      y \\
      z
    \end{bmatrix}
    & =
    \begin{bmatrix}
      0 \\
      0 \\
      0
    \end{bmatrix}
   \end{align*}
\column{0.4\textwidth}
  Solução
  \begin{align*}
    -3y   & = 0     \\
    -3x+y & = 0     \\
    x     & = y = 0 \\
    z     & = t \in \R
  \end{align*}
  Um autovetor
  \[
    \vec{v}_2
    = \vetortri{0}{0}{1}
  \]
\end{columns}
\end{resolution}

%------------------------------------------------------------------------------%
\begin{resolution}{Autoespaço de \(\lambda=5\)}
  \[
    E_5
    = \left\{
      \begin{bmatrix}
        0 \\
        0 \\
        t
      \end{bmatrix}
      \st t \in \R
    \right\}
    = \ger\left\{
      \begin{bmatrix}
        0 \\
        0 \\
        1
      \end{bmatrix}
    \right\}
  \]
\end{resolution}

%------------------------------------------------------------------------------%
\endinput



\begin{frame}{Exemplo 4: autovalores}
    Os autovalores são as raízes do polinômio característico
    \begin{align*}
        p(\lambda)
        &=\det(B-\lambda I)\\
        &=
        \begin{vmatrix}
            4-\lambda & 1 & 0\\
            1 & 4-\lambda & 0\\
            0 & 0 & 2-\lambda
        \end{vmatrix}\\
        &=(2-\lambda)
        \begin{vmatrix}
            4-\lambda & 1\\
            1 & 4-\lambda
        \end{vmatrix}\\
        &=(2-\lambda)
        \left[(4-\lambda)^2-1\right]\\
        &=(2-\lambda)(3-\lambda)(5-\lambda)
    \end{align*}
\end{frame}

\begin{frame}{Exemplo 4: autovalores}
    Resolvendo
    \[
        (2-\lambda)(3-\lambda)(5-\lambda)=0
    \]

    Obtemos
    \[
        \lambda_{1}=2
        \qquad
        \lambda_{2}=3
        \qquad
        \lambda_{3}=5
    \]
\end{frame}

\begin{frame}{Exemplo 4: autovetores para $\lambda=2$}
    Resolvemos o sistema
    \begin{align*}
        (B-2I)\vec{x}&=\vec{0}\\
        \begin{bmatrix}
            2 & 1 & 0\\
            1 & 2 & 0\\
            0 & 0 & 0
        \end{bmatrix}
        \begin{bmatrix}
            x\\y\\z
        \end{bmatrix}
        &=
        \begin{bmatrix}
            0\\0\\0
        \end{bmatrix}
    \end{align*}

    A forma escalonada do sistema fornece
    \begin{align*}
        2x+y&=0\\
        x+2y&=0\\
        z&\in\mathbb{R}
    \end{align*}

    Logo
    \begin{align*}
        x&=0\\
        y&=0
    \end{align*}

    Portanto
    \[
        E_{2}=
        \operatorname{span}
        \left\{\vetortri{0}{0}{1}\right\}
    \]
\end{frame}

\begin{frame}{Exemplo 4: autovetores para $\lambda=3$}
    Resolvemos o sistema
    \begin{align*}
        (B-3I)\vec{x}&=\vec{0}\\
        \begin{bmatrix}
            1 & 1 & 0\\
            1 & 1 & 0\\
            0 & 0 & -1
        \end{bmatrix}
        \begin{bmatrix}
            x\\y\\z
        \end{bmatrix}
        &=
        \begin{bmatrix}
            0\\0\\0
        \end{bmatrix}
    \end{align*}

    Obtemos
    \begin{align*}
        x+y&=0\\
        z&=0\\
        y&=-x
    \end{align*}

    Portanto
    \[
        E_{3}=
        \left\{
            \vetortri{t}{-t}{0}
            \mid t\in\mathbb{R}
        \right\}
        =
        \operatorname{span}
        \left\{\vetortri{1}{-1}{0}\right\}
    \]
\end{frame}

\begin{frame}{Exemplo 4: autovetores para $\lambda=5$}
    Resolvemos o sistema
    \begin{align*}
        (B-5I)\vec{x}&=\vec{0}\\
        \begin{bmatrix}
            -1 & 1 & 0\\
            1 & -1 & 0\\
            0 & 0 & -3
        \end{bmatrix}
        \begin{bmatrix}
            x\\y\\z
        \end{bmatrix}
        &=
        \begin{bmatrix}
            0\\0\\0
        \end{bmatrix}
    \end{align*}

    Obtemos
    \begin{align*}
        -x+y&=0\\
        z&=0\\
        y&=x
    \end{align*}

    Portanto
    \[
        E_{5}=
        \left\{
            \vetortri{t}{t}{0}
            \mid t\in\mathbb{R}
        \right\}
        =
        \operatorname{span}
        \left\{\vetortri{1}{1}{0}\right\}
    \]
\end{frame}

\begin{frame}{Exemplo 4: resposta}
    \[
        \lambda_{1}=2
        \qquad
        \lambda_{2}=3
        \qquad
        \lambda_{3}=5
    \]

    Os autoespaços são
    \begin{align*}
        E_{2}&=\operatorname{span}
        \left\{\vetortri{0}{0}{1}\right\}\\
        E_{3}&=\operatorname{span}
        \left\{\vetortri{1}{-1}{0}\right\}\\
        E_{5}&=\operatorname{span}
        \left\{\vetortri{1}{1}{0}\right\}
    \end{align*}

    Os autovetores são todos os vetores não nulos pertencentes aos respectivos autoespaços
\end{frame}